\hypertarget{appendix-f90cee1be7a4}{}
\section*{Optical focus and binocular conflict}
\section{Optical focus and binocular conflict}\label{optical-focus-and-binocular-conflict}

\textbf{Curriculum depth:} Introduced · \textbf{First seen:} XR01 · \textbf{Latest development:} XR01

\subsection{Definition and mechanism}\label{definition-and-mechanism}

Refraction, aperture and lens power shape retinal image formation. Accommodation changes focus; vergence aligns binocular targets. A fixed optical image plane can disagree with disparity-induced depth, while sharpness alone does not establish correct binocular appearance.

\subsection{Prerequisites}\label{prerequisites}

\begin{itemize}
\tightlist
\item
  \hyperref[appendix-215c2570199b]{Light spectra and radiometry} --- The light quantity and units must be established before its interaction or visual effect is modeled.
\end{itemize}

\subsection{Used by}\label{used-by}

\begin{itemize}
\tightlist
\item
  \hyperref[appendix-7d356c396d1a]{Visual angle and calibrated rays}
\item
  \hyperref[appendix-8b76b0efd27f]{PSF OTF and MTF}
\item
  \hyperref[appendix-8dc848f14052]{Stereoacuity and binocular foveation}
\end{itemize}

\subsection{Recall and next depth}\label{recall-and-next-depth}

Compute 0.5D display demand versus 4D near-object demand under the toy geometry. Pending: real optical calibration and stereo display design. Written coverage is not a claim of learner mastery or accepted numbered-curriculum progress.

\subsection{Source and navigation}\label{source-and-navigation}

\href{Continuous\%20Notes/University\%20-\%20Worlds,\%20Images\%20and\%20Inference\%20-\%20Graphics,\%20Vision\%20and\%20Photography\#Focus\%20and\%20stereo\%20can\%20disagree\%20even\%20when\%20the\%20images\%20are\%20sharp}{Expanded teaching explanation} · \href{}{Independent coverage (01-Foundations-I, appendix-fd090d982a15 § XR01 - Independent seeing and representation source)} · \href{}{Dependency Map (02-Foundations-II, appendix-4ad47bdd269a)}
